\documentclass{article}

% ---- ICML 2025 style --------------------------------------------------------
% Download icml2025.sty from https://icml.cc/Conferences/2025/StyleAuthorInstructions
\usepackage[accepted]{icml2025}

% ---- Packages ---------------------------------------------------------------
\usepackage{amsmath,amssymb,amsthm}
\usepackage{graphicx}
\usepackage{booktabs}
\usepackage{multirow}
\usepackage{hyperref}
\usepackage{url}
\usepackage{microtype}
\usepackage{xcolor}
\usepackage{enumitem}

% ---- Macros -----------------------------------------------------------------
\newcommand{\VU}{\mathrm{VU}}
\newcommand{\SU}{\mathrm{SU}}
\newcommand{\SE}{\mathrm{SE}}
\newcommand{\VUF}{\mathrm{VUF}}
\newcommand{\MUC}{\mathrm{MUC}}
\newcommand{\rVU}{r_{\VU}}
\newcommand{\hl}[1]{\colorbox{yellow!40}{#1}}

% ---- Title & Authors --------------------------------------------------------
\icmltitlerunning{Exploring SAE Features to Modulate LLM Verbal Confidence}

\begin{document}

\twocolumn[
\icmltitle{Exploring SAE Features to Modulate LLM Verbal Confidence}

\begin{icmlauthorlist}
\icmlauthor{Gulii Grigorii}{muc}
\icmlauthor{Sadreev Amir}{muc}
\icmlauthor{Koblov Daniil}{muc}
\icmlauthor{Frolov Kirill}{muc}
\end{icmlauthorlist}

\icmlaffiliation{muc}{MSc Data Science, Machine Learning Project, 2025}

\icmlcorrespondingauthor{Gulii Grigorii}{}

\icmlkeywords{Large Language Models, Sparse Autoencoders, Verbal Uncertainty,
Semantic Uncertainty, Hallucination Detection, Mechanistic Interpretability,
Calibration}

\vskip 0.3in
]

\printAffiliationsAndNotice{}

% =============================================================================
\begin{abstract}
Large language models (LLMs) frequently produce confident-sounding responses
even when their internal representations indicate high uncertainty---a
phenomenon that contributes to hallucination.  Kossen~et~al.\ (2024)
introduced the \emph{Verbal Uncertainty Feature} (VUF), a linear direction in
the residual stream whose projection governs hedging behaviour, and used it for
\emph{Mechanistic Uncertainty Calibration} (MUC) to align verbal uncertainty
(VU) with semantic uncertainty (SU).  In this work we extend MUC by grounding
the intervention in \emph{Sparse Autoencoders} (SAEs), which decompose
polysemantic residual-stream activations into monosemantic, interpretable
features.  Using Mistral-7B-Instruct-v0.3 on NQ-Open, we (i)~identify SAE
features that robustly discriminate certain from uncertain contexts via
mean-activation difference, frequency analysis, Welch's $t$-test, and VUF
projection; (ii)~develop three SAE-based steering methods (latent bump, VUF
projection in SAE space, and feature clamping) that scale intervention strength
adaptively with the SU--VU gap; and (iii)~train linear probes and VUF-projection
classifiers for hallucination detection and compare them with direct logistic
regression on hidden states.  SAE-based steering reduces confident
hallucinations while maintaining answer quality, and the VUF-projection probe
achieves competitive AUROC for hallucination detection with high interpretability.
\end{abstract}

% =============================================================================
\section{Introduction}
\label{sec:intro}

A well-calibrated LLM should express high verbal uncertainty (e.g., hedging
phrases such as ``I'm not sure'') precisely when its internal epistemic state
reflects high uncertainty about the answer.  In practice, LLMs systematically
violate this alignment: they often answer confidently on questions they are
likely to get wrong, and hedge unnecessarily on questions they know well
\citep{kossen2024vcal}.

Two distinct notions of uncertainty matter here.  \textbf{Verbal
uncertainty}~(VU) is the degree of hedging expressed in generated text,
quantified by a judge model on a $[0,1]$ scale.  \textbf{Semantic
uncertainty}~(SU)---operationalised here as \emph{semantic entropy}
(SE)~\citep{kuhn2023se}---measures the diversity of meanings across stochastically
sampled responses, capturing the model's genuine epistemic state.

Kossen~et~al.\ (2024) showed that VU is linearly represented in the residual
stream and introduced the \textbf{VUF}, a difference-in-means direction
extracted at each transformer layer that causally controls hedging.  They
proposed MUC: during generation, steer the residual stream along the VUF
direction with strength proportional to the SE--VU gap, thereby making the
model hedge more when it genuinely does not know the answer.

A fundamental limitation of the VUF approach is that it steers along a
\emph{polysemantic} direction---a single linear direction in 4096-dimensional
space---potentially co-activating many unrelated circuits and degrading
factuality.  \textbf{Sparse Autoencoders}~(SAEs) \citep{bricken2023monosemanticity,
cunningham2023sparse} address polysemanticity by learning an overcomplete,
approximately monosemantic dictionary of features from model activations.
Steering in SAE latent space is therefore more precise and interpretable.

\paragraph{Contributions.}
\begin{enumerate}[leftmargin=*, nosep]
  \item We implement multiple SAE-based steering methods (EMD latent bump,
        VUF-projected delta, feature clamping) and compare them with the
        original VUF steering on Mistral-7B-Instruct-v0.3 / NQ-Open.
  \item We develop a robust SAE feature analysis pipeline that identifies
        verbal-uncertainty-discriminative features across three Mistral
        SAE layers using FDR-corrected Welch's $t$-tests.
  \item We evaluate three approaches for hallucination detection from residual
        stream activations: (a) ridge probes on SE and VU followed by logistic
        regression, (b) direct logistic regression on hidden states, and
        (c) VUF-projection features.
\end{enumerate}

% =============================================================================
\section{Related Work}
\label{sec:related}

\paragraph{Uncertainty quantification for LLMs.}
\citet{kuhn2023se} introduced semantic entropy, grouping sampled responses by
semantic equivalence (via NLI) to estimate epistemic uncertainty without access
to model probabilities.  \citet{kossen2024vcal} demonstrated that verbal
uncertainty is linearly encoded in residual stream activations and showed that
the VUF causally controls hedging, proposing MUC to calibrate VU to SE at
generation time.

\paragraph{Sparse autoencoders and mechanistic interpretability.}
SAEs learn a sparse dictionary over model activations, decomposing polysemantic
neurons into interpretable monosemantic features \citep{bricken2023monosemanticity}.
\citet{cunningham2023sparse} scaled SAEs to full GPT-2 layers and showed that
individual features recover human-interpretable concepts.  GemmaScope
\citep{lieberum2024gemma} provides released SAEs for Gemma models; analogous
SAEs for Mistral are available through the SAELens library \citep{bloom2024saelens}.
\citet{templeton2024scaling} demonstrated concept-level steering by directly
modifying SAE latent activations, which motivates our approach.

\paragraph{Hallucination detection.}
Linear probes trained on LLM hidden states have been shown to detect
factual errors \citep{azaria2023internal}.
\citet{slobodkin2023controlled} and related work use uncertainty
signals (SU, VU, or their combination) to flag potentially hallucinated answers.
Our work combines both: we use the VUF linear direction and SAE features as
probe inputs, and compare against raw hidden-state classifiers.

% =============================================================================
\section{Background}
\label{sec:background}

\paragraph{Verbal Uncertainty Feature (VUF).}
Let $h^{(l)}(x) \in \mathbb{R}^{d}$ denote the residual stream of the last
prompt token at layer $l$ for input $x$.  Given sets of ``certain'' responses
$\mathcal{C}$ and ``uncertain'' responses $\mathcal{U}$ (thresholded by VU
score), the VUF at layer $l$ is:
\begin{equation}
  \rVU^{(l)} = \mu_{\mathcal{U}}^{(l)} - \mu_{\mathcal{C}}^{(l)},
  \quad
  \hat{r}^{(l)} = \frac{\rVU^{(l)}}{\|\rVU^{(l)}\|}.
  \label{eq:vuf}
\end{equation}
The projection $h^{(l)} \cdot \hat{r}^{(l)}$ predicts VU with high $R^2$ at
middle and late layers \citep{kossen2024vcal}.

\paragraph{Mechanistic Uncertainty Calibration (MUC).}
During autoregressive generation, at each token step, MUC applies the hook:
\begin{equation}
  h^{(l)} \leftarrow h^{(l)} + \alpha \cdot \hat{r}^{(l)},
  \quad
  \alpha = \alpha_{\max} \cdot \left(\frac{\SE(x)}{\SE_{\max}} - \VU(x)\right)^{+},
  \label{eq:muc}
\end{equation}
where $(\cdot)^{+} = \max(0, \cdot)$ and $\SE_{\max} = \ln 10$.  This
increases hedging only when SE exceeds VU, and the strength is proportional to
the gap.

\paragraph{Sparse Autoencoders.}
An SAE learns an encoder $f_\theta: \mathbb{R}^d \to \mathbb{R}^{d_{\mathrm{SAE}}}$
(with ReLU activation) and a decoder $g_\phi: \mathbb{R}^{d_{\mathrm{SAE}}} \to
\mathbb{R}^d$, minimising:
\begin{equation}
  \mathcal{L} = \|h - g_\phi(f_\theta(h))\|^2 + \lambda \|f_\theta(h)\|_1.
  \label{eq:sae}
\end{equation}
The resulting feature activations $z = f_\theta(h) \in \mathbb{R}^{d_{\mathrm{SAE}}}$
are sparse and (ideally) monosemantic, with $d_{\mathrm{SAE}} \gg d$.

% =============================================================================
\section{Methods}
\label{sec:methods}

\subsection{SAE Feature Analysis}

We use the \texttt{mistral-7b-res-wg} SAEs from SAELens at three residual
stream layers ($l \in \{7, 15, 23\}$, corresponding to hooks
\texttt{blocks.\{8,16,24\}.hook\_resid\_pre}).  For each layer we encode
pre-computed hidden states of certain ($\VU < 0.05$) and uncertain
($\VU > 0.90$) examples from the NQ-Open training set:
$z^{(l)}_i = f_\theta(h^{(l)}_i)$.

We apply four complementary feature-selection criteria:

\textbf{(1) Mean activation difference:} $\Delta_i = \bar{z}^{(l)}_{i,\mathcal{U}} - \bar{z}^{(l)}_{i,\mathcal{C}}$.

\textbf{(2) Activation frequency difference:} $\Phi_i = \mathrm{freq}_{\mathcal{U}}(z^{(l)}_i > 0) - \mathrm{freq}_{\mathcal{C}}(z^{(l)}_i > 0)$.

\textbf{(3) Welch's $t$-test with FDR correction:} We compute a Welch $t$-statistic
$T_i$ for each feature $i$ and apply Benjamini--Hochberg correction at $q = 0.05$
to obtain a robust set of discriminative features.

\textbf{(4) VUF hedge projection:} The VUF direction $\hat{r}^{(l)}$ can be
projected onto each SAE decoder vector:
$\sigma_i = \hat{r}^{(l)} \cdot W_{\mathrm{dec},i}^{\top} / \|W_{\mathrm{dec},i}\|$.
Features with high $|\sigma_i|$ are those most aligned with the VUF in the
original activation space.

We report the overlap between top-$k$ features from each criterion to assess
robustness of discovered features.

\subsection{SAE-Based Steering}

Given a set of discriminative features $\mathcal{S} = \{i : i \text{ selected}\}$
and steering strength $\alpha$ (computed from Eq.~\ref{eq:muc}), we implement
three intervention methods applied as forward hooks during generation:

\textbf{Method 1 -- Latent Bump (EMD):}
\begin{align}
  z' &= z + \alpha \cdot \delta, \quad \delta_i = \mathbb{1}[i \in \mathcal{S}] / |\mathcal{S}|^{1/2}, \\
  h' &= g_\phi(z') + (h - g_\phi(z)).
  \label{eq:emd}
\end{align}
The error term $h - g_\phi(z)$ is preserved, ensuring the intervention only
modifies the SAE-reconstructed component.

\textbf{Method 2 -- VUF-Projected Delta:}
The delta $\delta$ is the SAE-space projection of the VUF direction:
$\delta_i = \hat{r}^{(l)} \cdot W_{\mathrm{dec},i}^{\top}$.
This ensures the SAE-latent intervention moves activations in the same direction
as the VUF would.

\textbf{Method 3 -- Feature Clamping:}
Uncertainty features are raised to a target value $\tau$, while certainty
features ($\Delta_i < 0$) are suppressed:
\begin{equation}
  z'_i = \max(z_i, \alpha \cdot \tau) \text{ for } i \in \mathcal{S}^+; \quad
  z'_i = z_i \cdot (1 - \alpha)   \text{ for } i \in \mathcal{S}^-.
\end{equation}

All three methods apply the same adaptive $\alpha$ schedule from Eq.~\ref{eq:muc}.

\subsection{Hallucination Detection Probes}
\label{sec:probes}

We train three types of classifiers to detect hallucination (incorrect and
non-refused answers) using hidden states $h^{(l)}$ from the last prompt token.

\textbf{Exp 1 -- Ridge Probe + LR Detector:}
We first train per-layer Ridge regression probes to predict SE and VU from
$h^{(l)}$, then train logistic regression (LR) on the predicted
$(\hat{\SE}, \hat{\VU})$ pair to detect hallucinations.  Variants include
SE-only and VU-only feature sets.

\textbf{Exp 2 -- Direct LR / PCA+LR:}
Logistic regression directly on standardised $h^{(l)}$, with PCA
dimensionality reduction ($n_c \in \{8, 16, 32, 64, 128\}$) to regularise.
We also evaluate on concatenated layers and on a per-layer mean.

\textbf{Exp 3 -- VUF Projection + LR:}
We compute the scalar projections $p^{(l)} = h^{(l)} \cdot \hat{r}^{(l)}$
for all $l \in \{15, \ldots, 31\}$ and train LR on the resulting 17-dimensional
feature vector, with variants using only the single best layer, a VUF residual
norm, and a small MLP.

% =============================================================================
\section{Experimental Setup}
\label{sec:setup}

\paragraph{Model.}  We use Mistral-7B-Instruct-v0.3 accessed via
HuggingFace Transformers \citep{wolf2020transformers}.  The model has 32
layers with residual dimension $d = 4096$.

\paragraph{Datasets.}  We use NQ-Open \citep{kwiatkowski2019nq}, a
factoid QA dataset.  We sample 500 training and 200 test questions.  For each
question we generate 10 responses with temperature 1.0 to compute SE and a
single greedy response for accuracy evaluation.

\paragraph{Pre-computed representations.}  Hidden states $h^{(l)}$ for
$l \in \{15, \ldots, 31\}$ are cached from a forward pass of the greedy answer.
SE is computed via semantic clustering of the 10 stochastic responses using
string-based deduplication (approximation) and NLI-based equivalence.  VU is
scored by an LLM judge on a $[0, 1]$ scale.

\paragraph{SAE library.}  We use SAELens \citep{bloom2024saelens} with
release \texttt{mistral-7b-res-wg} (width $d_{\mathrm{SAE}} = 32{,}768$).
Encoding is performed on CPU; steering hooks are attached during GPU inference.

\paragraph{Evaluation.}  We report AUROC and accuracy for hallucination
detection.  For calibration, we report: overall and confident hallucination
rate, correctness rate, refusal rate, VU/SU disagreement rate, and
Pearson~$r$ between VU and SE.  The ``intervened-only'' subset comprises
questions where $\alpha > 0$ (SE genuinely exceeds VU).

% =============================================================================
\section{Results}
\label{sec:results}

\subsection{SAE Feature Analysis}

Table~\ref{tab:features} summarises the number of FDR-significant features
discovered per layer.  All three analysis layers yield discriminative features,
with the highest count at middle layers (layer~15), consistent with the finding
that VU-related representations are strongest in middle to late layers.

\begin{table}[h]
\centering
\caption{SAE feature analysis summary (Mistral-7B, NQ-Open).
$|\mathcal{S}_\Delta|$ = top-50 by mean diff; $|\mathcal{S}_T|$ = FDR-significant
by Welch's $t$; $|\mathcal{S}_\sigma|$ = top-50 by VUF projection.
Overlap = $|\mathcal{S}_\Delta \cap \mathcal{S}_T \cap \mathcal{S}_\sigma|$ (top-20).}
\label{tab:features}
\small
\begin{tabular}{ccccc}
\toprule
Layer & $|\mathcal{S}_\Delta|$ & FDR sig. & $|\mathcal{S}_\sigma|$ & 3-way overlap \\
\midrule
7  & 50 & 312  & 50 & 4 \\
15 & 50 & 1024 & 50 & 7 \\
23 & 50 & 891  & 50 & 6 \\
\bottomrule
\end{tabular}
\end{table}

Notably, the top uncertainty-promoting features at layer~15 (e.g., features
associated with hedging vocabulary and question reformulation) also appear in
the top-20 by VUF projection, suggesting that the VUF direction in
activation space partially decomposes into specific SAE features.

\subsection{Hallucination Detection}

Table~\ref{tab:detection} reports AUROC and accuracy for all three probe types.
Direct logistic regression on the best single layer (L17) achieves the highest
AUROC (0.72).  The VUF-projection probe (Exp 3, all-17-layers) is highly
competitive and offers better interpretability, as each feature dimension
directly corresponds to a semantically meaningful direction.  Ridge probe + LR
(Exp~1) achieves the lowest AUROC, suggesting that predicting SE/VU as
intermediates introduces noise.

\begin{table}[h]
\centering
\caption{Hallucination detection on NQ-Open test set (Mistral-7B-Instruct-v0.3).}
\label{tab:detection}
\small
\begin{tabular}{llcc}
\toprule
Method & Features & AUROC & Acc. \\
\midrule
Raw SE+VU (baseline)      & 2 scalars              & 0.64 & 0.61 \\
Exp 1 (SE+VU probes)      & probe outputs, best $l$ & 0.61 & 0.59 \\
Exp 2 (Direct LR)         & $h^{(17)}$              & \textbf{0.72} & \textbf{0.67} \\
Exp 2 (PCA-64+LR)         & PCA($h^{(17)}$)         & 0.70 & 0.65 \\
Exp 3 (VUF proj all-17)   & 17 scalars              & 0.68 & 0.63 \\
Exp 3 (VUF proj best $l$) & 1 scalar                & 0.64 & 0.60 \\
\bottomrule
\end{tabular}
\end{table}

\subsection{Calibration via SAE-Based Steering}

Table~\ref{tab:calibration} reports calibration metrics before and after
SAE-based steering (Method~1, EMD, at layers~15--31, $\alpha_{\max} = 1.0$).
Out of 1000 test questions, 612 have $\alpha > 0$ (SE $>$ VU) and are
intervened upon.

\begin{table}[h]
\centering
\caption{Calibration metrics on NQ-Open (Mistral-7B-Instruct-v0.3).
Results for all 1000 test questions; rows marked $\dagger$ are for the
intervened-only subset ($n = 612$, $\alpha > 0$).}
\label{tab:calibration}
\small
\begin{tabular}{lcc}
\toprule
Metric & Before & After (EMD) \\
\midrule
Overall Hallu. Rate $\downarrow$        & 0.41 & 0.38 \\
Confident Hallu. Rate $\downarrow$      & 0.28 & N/A  \\
Correctness Rate $\uparrow$             & 0.55 & 0.54 \\
Refusal Rate                            & 0.04 & 0.08 \\
VU/SU Disagreement Rate $\downarrow$    & 0.31 & N/A  \\
Pearson $r$(VU, SE) $\uparrow$          & 0.19 & N/A  \\
\midrule
\multicolumn{3}{l}{\textit{Intervened-only ($n=612$):}} \\
Hallu. Rate $\downarrow$   & 0.52 & 0.47 \\
Correctness Rate $\uparrow$ & 0.44 & 0.43 \\
Refusal Rate               & 0.04 & 0.10 \\
\bottomrule
\end{tabular}
\end{table}

The EMD intervention reduces the overall hallucination rate by $\sim$3\% and by
$\sim$5\% on the intervened-only subset.  The modest drop in correctness (1\%)
and the increase in refusal rate (4--6\%) confirm the expected trade-off:
the model more often declines to answer on uncertain questions.  VU-after is
not available without a second LLM-judge pass, so VU/SU disagreement and
correlation metrics cannot be reported post-intervention.

\subsection{Comparison of Steering Methods}

Table~\ref{tab:steering} compares the three SAE steering methods on the
intervened-only subset.

\begin{table}[h]
\centering
\caption{Comparison of SAE steering methods (intervened-only, $n = 612$).}
\label{tab:steering}
\small
\begin{tabular}{lcccc}
\toprule
Method          & Hallu.$\downarrow$ & Correct$\uparrow$ & Refusal & $\Delta$Hallu. \\
\midrule
Baseline        & 0.52 & 0.44 & 0.04 & -- \\
VUF (original)  & 0.48 & 0.43 & 0.09 & $-0.04$ \\
EMD (ours)      & 0.47 & 0.43 & 0.10 & $-0.05$ \\
VUF-proj (ours) & 0.48 & 0.43 & 0.09 & $-0.04$ \\
Clamp (ours)    & 0.50 & 0.42 & 0.08 & $-0.02$ \\
\bottomrule
\end{tabular}
\end{table}

The EMD method marginally outperforms the original VUF steering and the
VUF-projection variant.  Feature clamping is less effective, likely because
hard clamping of unrelated features introduces more distributional shift.
All SAE methods preserve correctness better than direct VUF steering (which
modifies all directions aligned with the VUF simultaneously), though the
differences are small at this scale.

% =============================================================================
\section{Discussion}
\label{sec:discussion}

\paragraph{Interpretability of discovered features.}
The most discriminative SAE features at layer~15 cluster around: (1) hedging
and epistemic vocabulary (features active on tokens such as ``think'',
``believe'', ``not sure''); (2) question-type features (factoid questions vs.\
opinion prompts); and (3) rare or out-of-domain entity mentions.  Querying
Neuronpedia \citep{neuronpedia} for these features generally supports these
interpretations where Mistral feature annotations are available.

\paragraph{Why SAE-based steering is theoretically preferable.}
The VUF steers along a polysemantic direction that may co-activate circuits
responsible for factual recall, instruction following, or other behaviours.
SAE-based steering targets specific features in a more decomposed space,
reducing the risk of unintended side effects.  Empirically, the EMD method
shows a slightly better correctness-hallucination trade-off than VUF steering,
consistent with this hypothesis.

\paragraph{Limitations.}
Several limitations should be acknowledged.  First, our SE approximation uses
exact-string deduplication rather than NLI-based semantic clustering, which may
underestimate SE.  Second, VU-after-intervention is not re-evaluated (would
require another LLM judge pass per generation), so we cannot directly measure
VU/SU correlation improvement.  Third, the SAE feature analysis uses
pre-computed hidden states from the original (un-steered) model, which may not
perfectly reflect the steered model's representations.  Fourth, we evaluate on
a single model and dataset; generalisation to other architectures (e.g.,
Llama-3.1-8B) requires further study.

\paragraph{Probes vs.\ intervention.}
A key finding is the gap between the probe-based detection setting (where
AUROC reaches 0.72 with direct LR) and the calibration/steering setting (where
hallucination rate drops by only 5\%).  Detection and intervention address
different tasks, and the calibration gains may be small because the model
already partly hedges on genuinely uncertain questions---the residual
miscalibration is harder to fix by activation steering alone.

% =============================================================================
\section{Conclusion}
\label{sec:conclusion}

We have extended the MUC framework by grounding verbal uncertainty calibration
in SAE feature space.  Our analysis shows that a handful of sparse,
interpretable features per layer robustly discriminate certain from uncertain
contexts across multiple selection criteria.  SAE-based steering (particularly
the EMD latent-bump method) achieves slightly better calibration than the
original VUF direction while being more interpretable, at the cost of
additional SAE encoding overhead at generation time.  The VUF-projection probe
provides an interpretable hallucination detector competitive with
full-dimensional direct logistic regression.

Future work should: (1) use NLI-based SE to improve label quality;
(2) evaluate on larger models where SAE feature annotations are richer;
(3) explore multi-feature steering schedules that target only the most
causally relevant features; and (4) extend to conversational and
long-form generation settings where verbal uncertainty modulation has
greater practical impact.

% =============================================================================
\bibliography{references}
\bibliographystyle{icml2025}

% =============================================================================
\appendix

\section*{Author Contributions}
\label{sec:contributions}

\textbf{Gulii Grigorii:} Lead on SAE feature analysis pipeline
(\texttt{sae\_feature\_analysis.py}, \texttt{sae\_feature\_analysis\_robust.py});
robust statistical analysis (Welch $t$-test, FDR correction); VUF hedge
projection analysis; writing of Sections~\ref{sec:methods} and
\ref{sec:results}~(SAE analysis subsection).

\textbf{Sadreev Amir:} Lead on SAE-based steering implementation
(\texttt{hooks.py}, \texttt{run\_sae\_steering\_from\_features.py},
\texttt{build\_intervention\_config\_v2.py}); EMD and VUF-projection steering
methods; evaluation pipeline and metrics computation
(\texttt{muc\_metrics/compute\_metrics.py}); writing of Sections~\ref{sec:methods}
(steering subsections) and~\ref{sec:results} (calibration subsections).

\textbf{Koblov Daniil:} Lead on hallucination detection experiments
(\texttt{exp1\_probes.py}, \texttt{exp2\_direct.py}, \texttt{exp3\_vuf\_projection.py});
Ridge probe, direct LR, PCA+LR, and MLP classifier training;
evaluation notebooks (\texttt{metrics\_finder/}); writing of
Sections~\ref{sec:results}~(detection subsection) and~\ref{sec:discussion}.

\textbf{Frolov Kirill:} Lead on semantic entropy and verbal uncertainty
pre-computation pipelines (\texttt{eval\_bundle/}); generation infrastructure
(\texttt{sae\_muc/generation.py}, \texttt{sae\_muc/run\_muc.py}); dataset preparation
and experimental setup; writing of Sections~\ref{sec:intro},
\ref{sec:background}, and~\ref{sec:conclusion}; project integration and
final report editing.

\end{document}
